\subsection{Quantized nonlinear filters and finite posterior labels}
\label{sec:filter-quantization-comparison}
The closest algorithmic comparison is with quantization of nonlinear
filters, not only with finite-window control. Pag\`es and Pham
\cite{PagesPham2005} develop marginal and Markovian quantizations of
the signal and a finite-dimensional filtering recursion, with error
bounds. Pag\`es and Sagna \cite[Section 6, Theorem 6.3 and
Remark 6.4]{PagesSagna2018} study a Markov signal with discrete-time
observations. Under their local regularity, moment and quantizer
hypotheses, the normalized-filter error is controlled by a weighted
sum of squared signal-quantization errors, with explicit likelihood
normalization factors. Their Remark 6.4 relates this estimate to the
globally Lipschitz case. Thus quantization followed by a stable
filtering recursion, including accumulation of preceding approximation
errors, is not a new principle of the present paper.

\paragraph{Object and resource.}
A signal grid with $K_n$ sites represents the approximate posterior by
weights on those sites; the grid size is not the number of possible
posterior weight vectors. In the present resource model, the entire
retained history-dependent state is one of $M_n$ labels. Each label
selects one fixed reachable prediction representative. The latter is
read-only numerical data, while its index is the charged resource.
Replacing a continuous weight vector by such a label requires a further
quantization of the attainable posterior image and a causal rule that
does not recover discarded history. Conversely, the present label law
is not a bound on total storage, offline construction, or arithmetic
workspace for a signal-grid implementation.

\paragraph{Observation model and norm.}
Our latent variable is fixed during an experiment, commands are observed
inputs, and reports come from a uniformly positive finite detector.
The query is selected independently after retention. The loss is the
weighted squared error of its physical success probability. Raw moments
preserve this norm up to fixed factors; an inverse covariance or an
inverse collision scale is not used. The monomial theorem fixes a
full-support prior, whereas the algebra theorem permits arbitrary priors
under bounded multiplication closure. These assumptions specify
subclasses, not an extension of every Markov-signal observation model.

\paragraph{Upper bounds and matching converses.}
An estimate for a supplied quantization scheme does not by itself
identify the infimum over all $M_n$-label posterior filters. Theorems
\ref{thm:resolution-main} and \ref{thm:algebra-multistep} give matching
bounds for that infimum and the corresponding checkpoint problems.
Their lower bounds use one unconditional command-and-report law and
allow independent coding randomness. Their upper bounds use reachable
representatives in one causal recursion. Neither the existence of such
a recursion nor a generic stability inequality identifies the attained
collision or covariance profile.

\paragraph{Degeneration and scope of the comparison.}
In the monomial class the constants remain uniform when additive
exponents collide, with the acquired-dimensional truncation intact.
In the algebra class they remain uniform through vanishing covariance
eigenvalues and prior support loss. The direct ellipsoid covering
argument in the latter case is elementary; its experiment-level content
is the whole-history product bound, the prior-uniform unconditional
acquisition mass, and the raw reachable-state update. We do not infer
from a filtering upper bound either these uniform geometric facts or a
matching posterior-label converse. Nor do we assert that the cited
filtering results are subsumed: their signal models, approximation
resources and norms differ. The theorem-level comparison here concerns
the displayed results, not an exhaustive priority claim.
